\section{Discussion and limitations}
\label{sec:discussion}

\paragraph{Interpretation.} On units where the task leaves a behaviour open, a hidden-test ``detection'' is
agreement with the benchmark's behaviour, a convention two generator models' task-only solutions also follow (89\% on
held-out units): auditors flag the program that departs from it more often (Section~\ref{sec:e6}). We cannot
separate a natural reading of the task from a convention the models share with the benchmark.
Adding one sentence that states the requirement turns that agreement into verified diagnosis; we could not resolve
whether a larger reasoning budget changes the result once it is stated. Natural failures are observational and
classified by models, but point in the same direction: both classifiers agree that the task's words settle fewer than
a quarter of failures, and the generators' other solutions often behave like the ``defective'' program
(Section~\ref{sec:natural}). The earlier scientific studies are consistent with the distinction without testing it
(Section~\ref{sec:transfer}): naming the deliverable cut false flags on correct scientific code from 69\% to 36\%,
and where code and inputs determine a reported value, recomputation catches fabrications that consistency checks
cannot.

\paragraph{Practical consequences.} Alongside recall, an audit benchmark should report which of its cases the
task determines. It should not count agreement with an unstated convention as detection. An auditor that is
unsure whether a behaviour is required should be scored for asking, not only for flagging or passing.
In deployment, an auditor's value on such behaviour depends on whether the requirement is written down.
Discovering requirements that nobody wrote down, which we did not measure, is the harder problem.

\paragraph{Limitations.}\label{sec:limits} \emph{Constructed units.} The randomised-contract units are execution-validated;
two models, rather than people, judge them admissible. Both models also implement the rules, so construction and
judging are not independent. The judges outside construction (E9) are the auditor models.
We make no claim about the frequency of such ambiguity. The 41 held-out units, selected from 406 problems,
are concentrated in MBPP+ (37). Task-only solutions include followers of both rules on only two units.
\emph{Benchmarks.} HumanEval+ and MBPP+ supply the natural failures and 41 of the 59 held-out randomised-contract units;
the 18 BigCodeBench units are reference edits (Section~\ref{sec:transfer}). The convention advantage remains
untested on unedited library code, where more behaviour is unstated. \emph{Auditors.} We tested two vendors' models at two reasoning budgets with one rulebook.
Other auditors, agent scaffolds or execution-based test generation may sit closer to the ceiling before the
reveal, although Lemma~\ref{lem:symmetric} bounds what any of them can do on the unrevealed arms.
\emph{Verification.} To verify a diagnosis, execution must confirm a witness's disagreement with the contract's oracle
using the benchmark's own comparator. The auditor's stated expected value must also be correct.
We exclude diagnoses stated only in prose, which can understate performance for auditors with correct but
unexecutable findings. The frozen scorer also gives no credit when the oracle raises. Crediting, post hoc, the two units whose
rule A requires an exception raises E6-5 from 76.8 to 79.3 and changes no verdict (Appendix~\ref{app:e6}).
\emph{Natural failures.} Models classify the natural failures (Cohen's $\kappa$ 0.71 on held-out items);
the behavioural check reuses the generators' own solutions.
\emph{Data.} The data study injects one fault per item into twelve public tables; its conclusion is specific to
faults that a clean reference sample reveals. \emph{Discovering requirements.}
In a registered follow-up (E8, Appendix~\ref{app:e6}), an auditor asked to list unspecified behaviours named the open
behaviour in 97.6\% of held-out readings but also the decided behaviour of positive controls in 80.8\%.
Selectivity was 16.8 points [10.5, 23.2], not resolved against the registered 20-point threshold. \emph{Amendments} to E6's infrastructure rules were registered before any held-out analysis (Appendix~\ref{app:e6}).
